\section{Introducción}
\label{sec:introduccion}

Los gestos de la mano son una forma natural de dar órdenes a una máquina sin mando ni contacto.
Esta ventaja es especialmente útil en entornos ruidosos o industriales, donde la voz no es
fiable, o cuando el operador tiene las manos ocupadas y no puede manejar un mando
\citep{pisharady2015review}. Para un robot, sin embargo, no basta con reconocer gestos: hace falta
un \emph{vocabulario de órdenes} con un significado preciso para cada uno y, sobre todo, una forma
segura de no hacer nada. Por eso este trabajo plantea el problema como la clasificación de un
conjunto cerrado de órdenes más una \emph{clase nula} (\cls{no\_gesture}) que deja el robot
quieto, en lugar de como un reconocimiento de gestos arbitrarios.

Este planteamiento impone tres requisitos que guían todo el diseño experimental:
\begin{enumerate}
  \item \textbf{Robustez ante personas nuevas.} El robot tendrá que obedecer a personas que el
    modelo nunca vio. Por tanto, la estimación del rendimiento debe hacerse con personas ausentes
    del entrenamiento; una partición aleatoria por imagen sería optimista, porque la misma persona,
    con el mismo fondo y la misma iluminación, aparecería a ambos lados \citep{kaufman2012leakage}.
  \item \textbf{Una clase nula fiable.} Confundir una mano en reposo con una orden mueve el robot
    sin querer, que es el error más costoso en la práctica.
  \item \textbf{Bajo coste de inferencia.} La clasificación debe ejecutarse en tiempo real en un
    ordenador portátil, lo que favorece modelos pequeños.
\end{enumerate}

\paragraph{Objetivos.}
El objetivo principal es alcanzar al menos un 85\,\% de exactitud en test con personas no vistas,
usando una red convolucional diseñada y entrenada desde cero en PyTorch \citep{paszke2019pytorch}.
La red debe incluir convolución, \emph{pooling}, normalización por lotes y \emph{dropout}, y su
entrenamiento debe usar parada anticipada sobre la pérdida de validación con restauración de los
mejores pesos. Además, nos propusimos analizar el modelo con una matriz de confusión, mapas
Grad-CAM y un estudio de ablación (con y sin \emph{dropout}, normalización por lotes y aumento
de datos). Como extensiones opcionales, el plan incluía comparar con un ResNet18 preentrenado y
construir una demo en tiempo real que controle un robot simulado en ROS\,2.

\paragraph{Contribuciones.}
\begin{itemize}
  \item Un \emph{pipeline} de datos reproducible sobre HaGRID: recorte de manos a partir de las
    cajas anotadas, una clase nula construida con las manos que no hacen el gesto y una partición
    estratificada \emph{por persona} sin personas compartidas entre particiones
    (\cref{sec:datos}).
  \item GestureCNN, una red de \ParamsFinal{} parámetros que alcanza un \AccTestFinal{} de
    exactitud en test, con un intervalo de confianza que tiene en cuenta la dependencia entre
    imágenes de una misma persona (\cref{sec:resultados}).
  \item Un estudio de ablación con dos semillas y pruebas de McNemar pareadas, que separa los
    efectos reales de la variación entre ejecuciones (\cref{sec:ablacion}).
  \item Una comparación con ResNet18 que incluye el coste computacional (parámetros, operaciones y
    latencia) y un análisis cualitativo de los errores con Grad-CAM.
\end{itemize}

El resto del documento se organiza así: la \cref{sec:relacionado} resume el trabajo relacionado;
la \cref{sec:datos} describe los datos y la partición; la \cref{sec:metodos}, el preprocesado, la
arquitectura y el protocolo experimental; la \cref{sec:resultados} presenta los resultados; la
\cref{sec:demo}, el diseño de la demo; la \cref{sec:discusion} discute los resultados y sus
limitaciones, y la \cref{sec:conclusiones} concluye.
